%! TeX root = ../tesi.tex
% =============================================================================
% capitolo5_conclusioni.tex
% Capitolo 5: Conclusioni e sviluppi futuri
%
% LINEE GUIDA VITALI:
%  L'ultimo capitolo chiude il cerchio con l'Introduzione.
%
%  A) RIPRESA DELLA TESI INIZIALE:
%     Ricollega tutto all'affermazione centrale dell'Introduzione.
%     "In questa tesi abbiamo dimostrato che [claim]..."
%     Sintetizza come ogni capitolo ha contribuito alla dimostrazione.
%
%  B) PUNTI DI ORGOGLIO (contributi originali):
%     Cosa hai fatto che non era stato fatto prima?
%     Sii specifico: "Il principale contributo originale e'..."
%     Non essere modesto oltre misura, ma nemmeno esagerare.
%
%  C) PUNTI DI MODESTIA (limiti riconosciuti):
%     Cosa NON fa il sistema? Quali assunzioni hai fatto?
%     Questo NON e' un punto di debolezza: dimostra maturita' scientifica.
%
%  D) SVILUPPI FUTURI:
%     3-5 direzioni concrete e motivate (non vaghe).
%     Esempio: "Un primo sviluppo futuro potrebbe essere l'integrazione con X
%               per risolvere Y, che attualmente e' un limite del sistema."
%
%  LUNGHEZZA SUGGERITA: 3-5 pagine.
% =============================================================================

% --- A. RIPRESA DELLA TESI ---
% "In questa tesi abbiamo dimostrato che..." Ricollega al claim dell'Introduzione.

% --- B. PUNTI DI ORGOGLIO ---
% I contributi originali del lavoro, in modo diretto e preciso.

% --- C. PUNTI DI MODESTIA ---
% Limiti del lavoro: cosa non fa, cosa hai assunto, margini di errore.

% --- D. SVILUPPI FUTURI ---
% 3-5 direzioni future concrete e motivate.

Scrivi qui il contenuto del capitolo 5 sulle conclusioni e gli sviluppi futuri.
